\documentclass[10pt,letterpaper]{article}
\usepackage{cornell-notes}

\title{Clause 28.05.03.04: Random Number Engine Requirements}
\author{}
\date{}
\setDocTitle{Clause 28.05.03.04: Random Number Engine Requirements}
\setDocSubtitle{C++ 2024}
\setDocOwner{C++ 2024 Cornell Notes}
\setDocDate{\today}
\setCornellCollection{C++ 2024 Cornell Notes}
\setCornellUnitType{Clause}
\setCornellUnitNumber{28.05.03.04}
\setCornellUnitTitle{Random Number Engine Requirements}
\hypersetup{pdftitle={Clause 28.05.03.04: Random Number Engine Requirements},pdfsubject={Cornell notes},pdfauthor={}}

\begin{document}
\maketitle
\begin{CornellOverviewBox}{Overview}
\textbf{Classification:} Standard library - Normative\\[2pt]
\textbf{Learning objective:} Understand the rules, terminology, and semantic consequences associated with Random number engine requirements.
\end{CornellOverviewBox}\begin{CornellSummaryBox}{Summary}
{\sffamily\bfseries\color{Accent} Summary}\par\vspace{3pt}Section 28.5.3.4 focuses on Random number engine requirements. Apply the specified interface together with its mandates, effects, result conditions, exception behavior, and complexity constraints. When applying it, preserve the distinction between mandatory requirements, recommendations, permissions, and behavior deliberately left undefined, unspecified, or implementation-defined.
\end{CornellSummaryBox}

\end{document}
